\documentclass[11pt,a4paper]{article}
\usepackage[utf8]{inputenc}
\usepackage[T1]{fontenc}
\usepackage{lmodern}
\usepackage[italian]{babel}
\usepackage[margin=2.2cm]{geometry}
\usepackage{amsmath,amssymb}
\usepackage{booktabs,longtable,array}
\usepackage{enumitem}
\usepackage[hidelinks]{hyperref}
\setlist{itemsep=2pt,topsep=3pt}
\setlength{\parindent}{0pt}
\setlength{\parskip}{5pt}
\setlength{\emergencystretch}{3em}
\sloppy
\newcommand{\code}[1]{\texttt{#1}}
\newcommand{\Dom}[1]{\par\medskip\noindent\textbf{Domanda.} #1\par\nopagebreak}
\newcommand{\Ris}[1]{\noindent\textbf{Risposta.}\par\nopagebreak #1}
\newcommand{\Trap}[1]{\par\smallskip\noindent\textbf{Errore tipico all'orale.} #1\par}

\title{Domande d'orale sull'intero corso\\Statistics and Data Analysis (Week 1--4, L9--L20)}
\author{Giacomo Moretti}
\date{Ottobre 2026}

\begin{document}
\maketitle

\section*{Come usare questo documento}
Sono le domande in stile esame fatte in sessione di studio, con la risposta modello, in ordine di programma.
Le sorgenti sono le slide del corso (\code{Statistics\_course/00\_Reference/All\_slides\_merged.pdf}, lezioni Week~1--4 e L9--L20) e Cowan, \emph{Statistical Data Analysis}.
Le sezioni sulla Week~1 (probabilit\`a, CLT, lognormale, propagazione degli errori) e sui limiti superiori sono scritte a partire dal programma del corso e da Cowan, perch\'e nelle slide non ci sono: da confrontare con i notebook (\code{Week1\_probability.ipynb}, \code{Week2\_CountingExperiment.ipynb}).
Il collegamento al progetto (punto (d) della sezione su unfolding e OmniFold) \`e da completare dopo aver riletto il report.

\begin{center}
\begin{tabular}{p{1.8cm}p{5.2cm}p{7cm}}
\toprule
Blocco & Lezioni & Temi \\
\midrule
1 & Week 1 & probabilit\`a, PDF/CDF, momenti, CLT, propagazione errori, Monte Carlo \\
2 & Week 2 & test di ipotesi, Neyman--Pearson, $p$-value \\
3 & Week 2--3 & stimatori, MLE, informazione di Fisher, MVB \\
4 & Week 3 & intervalli di confidenza, Wilks, profile likelihood \\
5 & Week 2--4 & limiti superiori, goodness of fit, combinazioni, intervalli credibili \\
6 & Week 4 & Bayes, MCMC, Bayes factor \\
7 & L9, L10 & bias--varianza, loss come NLL, Ridge/LASSO, CV \\
8 & L10, L11 & PCA, K-means, alberi, bagging/boosting \\
9 & L12, L13 & reti neurali, backprop, ottimizzatori, sbilanciamento \\
10 & L14, L13 & immagini, filtri, preparazione dei dati \\
11 & L14--L16 & CNN, RNN/LSTM, GNN \\
12 & L17, L18 & autoencoder, VAE, incertezza in ML \\
13 & L19, L20 & unfolding, OmniFold, progetto \\
14 & tutto & confronti e domande trabocchetto \\
\bottomrule
\end{tabular}
\end{center}

\clearpage
%==============================================================
\section{Probabilit\`a, momenti, CLT, propagazione degli errori e Monte Carlo (Week 1)}

\Dom{(a)~Quali sono gli assiomi di Kolmogorov? Scrivi la probabilit\`a condizionata, il teorema di Bayes e la legge delle probabilit\`a totali. Un test medico ha sensibilit\`a alta ma la malattia \`e rara: perch\'e un positivo non implica che sei malato? Cosa cambia tra interpretazione frequentista e bayesiana?
(b)~Come sono legate PDF e CDF? Come si definisce un quantile? Cosa sono densit\`a congiunte, marginali e condizionate, e quando due variabili sono indipendenti?
(c)~Definisci valore atteso, varianza, covarianza e coefficiente di correlazione. Indipendenti implica incorrelate? E il contrario? Come si trasforma una PDF sotto un cambio di variabile?
(d)~Enuncia il teorema del limite centrale e quando non vale. Perch\'e il prodotto di molte variabili positive d\`a una lognormale?
(e)~Come si propagano linearmente gli errori? Quando smette di valere?
(f)~Spiega hit-or-miss, rejection sampling e transformation method. A cosa serve il Monte Carlo in un'analisi?}

\Ris{
\textbf{(a)} Assiomi: (1)~$P(A)\ge0$; (2)~$P(S)=1$; (3)~se $A\cap B=\emptyset$ allora $P(A\cup B)=P(A)+P(B)$. Ne seguono $P(\bar A)=1-P(A)$, $P(\emptyset)=0$, $A\subset B\Rightarrow P(A)\le P(B)$ e $P(A\cup B)=P(A)+P(B)-P(A\cap B)$.
Condizionata: $P(A\mid B)=P(A\cap B)/P(B)$. $A$ e $B$ sono indipendenti se $P(A\cap B)=P(A)P(B)$, e allora $P(A\mid B)=P(A)$.
Probabilit\`a totali, con $S$ diviso in sottoinsiemi disgiunti $A_i$: $P(B)=\sum_i P(B\mid A_i)P(A_i)$. Bayes:
\[
P(A\mid B)=\frac{P(B\mid A)P(A)}{\sum_i P(B\mid A_i)P(A_i)} .
\]
Test medico: $P(D\mid +)=\dfrac{P(+\mid D)\pi_D}{P(+\mid D)\pi_D+P(+\mid\bar D)(1-\pi_D)}$. Se la prevalenza $\pi_D$ \`e piccola, anche con sensibilit\`a alta i falsi positivi, che vengono dalla grande popolazione dei sani, possono dominare: il posterior pu\`o essere basso.
Il teorema di Bayes \`e valido in qualsiasi interpretazione; cambia il significato del risultato.
Frequentista: $P(A)=\lim_{n\to\infty}(\text{volte in cui esce }A)/n$, per esiti di esperimenti ripetibili; le probabilit\`a si associano ai dati, non alle ipotesi (non ha senso $P(\text{supersimmetria})$).
Bayesiana (soggettiva): $P(A)$ \`e il grado di credenza che $A$ sia vera; si pu\`o assegnare una probabilit\`a a un'ipotesi, ma serve un prior. Nel test medico: per te ``ho il $X\%$ di probabilit\`a di essere malato'' (Bayes), per il medico ``il $X\%$ di chi risulta positivo \`e malato'' (frequentista).
I due approcci rispondono a domande diverse e coincidono con molti dati e prior deboli; la likelihood \`e il principio unificante.

\textbf{(b)} $F(x)=\int_{-\infty}^x f(x')dx'$ e $f=dF/dx$; $F$ \`e crescente, con $F(-\infty)=0$ e $F(\infty)=1$. Il quantile $x_\alpha$ \`e il valore con $F(x_\alpha)=\alpha$, cio\`e $x_\alpha=F^{-1}(\alpha)$; la mediana \`e $x_{0.5}$.
Densit\`a congiunta $f(x,y)$; marginale $f_x(x)=\int f(x,y)\,dy$; condizionata $f(y\mid x)=f(x,y)/f_x(x)$. Indipendenza: $f(x,y)=f_x(x)f_y(y)$ (e quindi la likelihood di dati iid \`e un prodotto).

\textbf{(c)} $E[x]=\int x f(x)dx$; $V[x]=E[(x-\mu)^2]=E[x^2]-\mu^2$; $\mathrm{cov}[x,y]=E[xy]-\mu_x\mu_y$; $\rho=\mathrm{cov}[x,y]/(\sigma_x\sigma_y)\in[-1,1]$.
Indipendenti $\Rightarrow$ incorrelate, ma non il contrario (esempio: $x$ simmetrico e $y=x^2$ hanno correlazione nulla ma sono dipendenti); $\rho$ misura solo la dipendenza lineare.
Cambio di variabile $y=y(x)$ monotono: $g(y)=f(x(y))\,\bigl|dx/dy\bigr|$; in pi\`u dimensioni il determinante dello jacobiano.
Esempio dal corso: un prior piatto in $\theta$ non \`e piatto in $\lambda=1/\theta$.

\textbf{(d)} Se $x_1,\dots,x_n$ sono iid con media $\mu$ e varianza finita $\sigma^2$, la somma tende a una Gaussiana con media $n\mu$ e varianza $n\sigma^2$; la media campionaria ha varianza $\sigma^2/n$. \`E il motivo per cui gli errori sono spesso gaussiani.
Non vale se la varianza \`e infinita (esempio: la Cauchy, la cui media campionaria non si restringe) o se le variabili sono fortemente correlate; la convergenza \`e lenta per distribuzioni molto asimmetriche.
Lognormale: se $y=\prod_i x_i$ con $x_i>0$, allora $\ln y=\sum\ln x_i$ \`e (per il CLT) gaussiano. Quindi $y$ \`e lognormale: asimmetrica, con coda destra pesante, positiva, con mediana $e^{\mu}$ e media $e^{\mu+\sigma^2/2}$.

\textbf{(e)} Per $y=f(x_1,\dots,x_n)$ linearizzando attorno ai valori medi:
\[
\sigma_y^2=\sum_{i,j}\frac{\partial f}{\partial x_i}\frac{\partial f}{\partial x_j}V_{ij}
=\nabla f^T\,V\,\nabla f ,
\]
dove $V$ \`e la matrice di covarianza. Per variabili incorrelate: somma in quadratura dei contributi $(\partial f/\partial x_i)^2\sigma_i^2$. Per somme si sommano in quadratura gli errori assoluti, per prodotti e rapporti quelli relativi; con le correlazioni il segno del termine misto conta.
Vale se $f$ \`e circa lineare entro l'intervallo di variazione degli $x_i$ (errori piccoli rispetto alla scala di variazione di $f$). Altrimenti si propaga con il Monte Carlo: si campionano gli $x_i$, si calcola $y$ e se ne guarda la distribuzione.

\textbf{(f)} Tre metodi per campionare una PDF $f(x)$ a partire da numeri uniformi:
\begin{itemize}
\item \textbf{Hit-or-miss}: si prende $x_i$ uniforme sul dominio e $y_i$ uniforme in $[0,M]$ con $M\ge\max f$; si tiene $x_i$ solo se $y_i<f(x_i)$.
\item \textbf{Rejection sampling}: come sopra ma si propone $x_i$ da una PDF $g(x)$ simile a $f$, con $M g(x)\ge f(x)$; $y_i$ uniforme in $[0,Mg(x_i)]$, si tiene se $y_i<f(x_i)$. L'efficienza \`e circa $1/M$.
\item \textbf{Transformation method}: si impone $r=F(x(r))$, quindi $x(r)=F^{-1}(r)$ con $r$ uniforme in $[0,1]$. Richiede di integrare $f$ e invertire la CDF.
\end{itemize}
Un computer usa PRNG (numpy: PCG64; ROOT: TRandom3/Mersenne Twister): algoritmi deterministici, con un seed che garantisce la riproducibilit\`a.
Usi del MC: stimare valori attesi e probabilit\`a (frazione di campioni), propagare errori non lineari, simulare esperimenti ripetuti (pseudo-esperimenti), studiare bias e varianza degli stimatori, ottenere la distribuzione delle statistiche di test, calcolare posterior (MCMC), addestrare algoritmi di ML.
Distribuzioni di riferimento: Poisson come limite della binomiale (e circa gaussiana per media grande); $\chi^2_n$ come somma di $n$ gaussiane standard al quadrato (media $n$, varianza $2n$); esponenziale per i tempi di decadimento.
}
\Trap{Dire che incorrelato significa indipendente; applicare il CLT senza varianza finita; confondere media e mediana della lognormale.}

\section{Test di ipotesi, Neyman--Pearson e $p$-value (Week 2)}

\Dom{Cerchi un segnale: $H_0$ = solo fondo, $H_1$ = segnale + fondo, e hai un classificatore con output $t(x)$.
(a)~Cosa dice il lemma di Neyman--Pearson e perch\'e il rapporto di verosimiglianza \`e ottimale?
(b)~Se NP d\`a gi\`a il test migliore, perch\'e si addestrano classificatori ML?
(c)~Il $p$-value risulta $0.003$: cosa significa e cosa non significa? Come \`e distribuito sotto $H_0$? Perch\'e in fisica delle particelle si chiede $Z>5$?}

\Ris{
\textbf{(a)} Per un test di dimensione $\alpha$ fissata, il test con la massima potenza contro $H_1$ usa la statistica
\[
t(x)=\frac{P(x\mid H_1)}{P(x\mid H_0)},
\]
con regione critica $t>t_c$ scelto in modo che $P(t>t_c\mid H_0)=\alpha$. Qualsiasi funzione monotona di $t$ (per esempio il logaritmo) d\`a lo stesso test.
Intuizione: si mettono nella regione critica i punti in cui $H_1$ \`e pi\`u probabile rispetto a $H_0$; a parit\`a di probabilit\`a spesa sotto $H_0$ ($\alpha$) si guadagna la massima probabilit\`a sotto $H_1$ (la potenza).

\textbf{(b)} NP richiede le PDF multivariate di $H_0$ e $H_1$, che di solito non abbiamo. Spesso c'\`e solo un MC, che d\`a campioni ma non una formula per $f(x\mid H)$: la simulazione passa per variabili latenti e l'integrale su di esse \`e intrattabile.
Stimare una PDF in alta dimensione \`e praticamente impossibile (curse of dimensionality).
Un classificatore addestrato con la BCE approssima il rapporto di verosimiglianza: trasforma un problema di stima di densit\`a (difficile) in uno di classificazione (facile). Lo stesso principio rende possibile OmniFold.

\textbf{(c)} Il $p$-value \`e la probabilit\`a, assumendo $H_0$, di osservare dati almeno altrettanto discrepanti da $H_0$:
\[
p=\int_{t_{\rm obs}}^{\infty} f(t\mid H_0)\,dt .
\]
Non \`e $P(H_0\mid\text{dati})$: in ambito frequentista $P(H_0)$ non esiste; servirebbe Bayes con un prior.
Sotto $H_0$ il $p$-value \`e uniforme in $[0,1]$; sotto $H_1$ tende a stare vicino a $0$.
$p=0.003$ significa: se $H_0$ fosse vera, un risultato cos\`i estremo capiterebbe lo $0.3\%$ delle volte.
In fisica delle particelle si chiede $Z>5$ ($p\approx3\cdot10^{-7}$) per la look-elsewhere effect (molti test, $p$-hacking), per sistematiche non previste e per evitare falsi annunci di scoperta.
}
\Trap{Dimenticare che il $p$-value dipende implicitamente dall'ipotesi alternativa, che definisce cosa significa ``pi\`u estremo''; confonderlo con $P(H_0)$.}

\section{Stimatori, MLE, Fisher e MVB (Week 2--3)}

\Dom{(a)~Scrivi la likelihood di $n$ misure indipendenti da $f(t;\tau)=\tfrac1\tau e^{-t/\tau}$ e ricava l'MLE di $\tau$.
(b)~L'MLE \`e sempre non distorto? Che propriet\`a ha asintoticamente?
(c)~Come stimi l'incertezza su $\hat\tau$ in almeno tre modi? Cosa dice la Minimum Variance Bound e quando l'MLE la raggiunge?}

\Ris{
\textbf{(a)} $L(\tau)=\prod_i \tfrac1\tau e^{-t_i/\tau}$, quindi $\ln L=-n\ln\tau-\tfrac1\tau\sum_i t_i$. Ponendo $\partial\ln L/\partial\tau=0$ si ottiene $\hat\tau=\tfrac1n\sum_i t_i$, la media campionaria. Qui \`e non distorto, con $V[\hat\tau]=\tau^2/n$.

\textbf{(b)} No. Esempio: per una Gaussiana lo stimatore $\hat\sigma^2=\tfrac1n\sum(x_i-\bar x)^2$ \`e distorto; quello non distorto ha $1/(n-1)$.
Inoltre l'MLE ha invarianza funzionale (l'MLE di $g(\theta)$ \`e $g(\hat\theta)$), e questo di solito rende distorti gli stimatori trasformati.
Asintoticamente ($n\to\infty$) \`e consistente, non distorto, gaussiano ed efficiente (raggiunge la MVB).

\textbf{(c)} Metodi:
\begin{itemize}
\item analitico (raramente praticabile);
\item MC: pseudo-esperimenti generati con $\hat\tau$ al posto del vero $\tau$, ricalcolo di $\hat\tau$, deviazione standard;
\item bootstrap: ricampionamento con reinserimento; funziona male con pochi dati e con parametri legati alle code;
\item seconda derivata: $\sigma^2=\bigl(-\partial^2\ln L/\partial\theta^2\bigr)^{-1}_{\theta=\hat\theta}$ (Minuit/Hesse);
\item grafico: l'intervallo con $2\Delta\mathrm{NLL}=1$, che permette errori asimmetrici (Minos).
\end{itemize}
Information inequality (MVB):
\[
V[\hat\theta]\ \ge\ \frac{\bigl(1+\partial b/\partial\theta\bigr)^2}{I(\theta)},\qquad I(\theta)=-E\!\left[\frac{\partial^2\ln L}{\partial\theta^2}\right],
\]
con $b$ il bias e $I$ l'informazione di Fisher; per stimatori non distorti $V\ge1/I$.
L'MLE raggiunge la MVB asintoticamente; a $n$ finito solo in casi particolari (come questo).
In pi\`u parametri $V\simeq I^{-1}$: la covarianza \`e l'inversa della matrice di Fisher.
}

\section{Intervalli di confidenza, Wilks e profile likelihood (Week 3)}

\Dom{Hai un parametro di interesse $\mu$ (eventi di segnale) e un parametro di disturbo $\nu$ (fondo).
(a)~Perch\'e ``il $68\%$ di probabilit\`a che $\mu$ sia nell'intervallo'' \`e sbagliato per un intervallo di confidenza? Come si costruisce con la Neyman construction e cosa significa coverage?
(b)~Cosa dice il teorema di Wilks? Con un POI e $5$ parametri di disturbo quanti gradi di libert\`a ha il $\chi^2$? Quale $2\Delta\mathrm{NLL}$ d\`a il $68\%$ con $1$ e con $2$ d.o.f.?
(c)~Perch\'e il frequentista fa il profiling di $\nu$ e il bayesiano marginalizza?}

\Ris{
\textbf{(a)} Per un frequentista $\mu$ \`e fisso e fluttua l'intervallo, perch\'e dipende dai dati. Il $68\%$ \`e la frazione di esperimenti ripetuti i cui intervalli contengono il vero $\mu$ (\emph{coverage}).
Neyman construction: si scandiscono i valori di $\mu$; per ciascuno si fa un test di ipotesi e si calcola il $p$-value dei dati osservati; l'intervallo \`e l'insieme dei $\mu$ con $p_\mu\ge\alpha$. Il vero valore \`e rifiutato con probabilit\`a $\alpha$, quindi la copertura \`e $1-\alpha$.
La frase giusta per l'intervallo credibile bayesiano \`e invece ``dati questi dati e questo prior, $\mu$ sta qui con il $68\%$''.

\textbf{(b)} Con $t_\theta=2\Delta\mathrm{NLL}(\theta)=-2\ln\bigl[L(\theta)/L(\hat\theta)\bigr]$, nel limite di grandi campioni e con condizioni di regolarit\`a $t_\theta\sim\chi^2_n$, con $n$ il numero di parametri su cui si fa il test.
Con un POI e $5$ nuisance profilati: $n=1$ (contano i POI, non i parametri totali).
Soglie al $68.3\%$: $2\Delta\mathrm{NLL}=1$ con $1$ d.o.f. (classico $\Delta\ln L=0.5$) e $2.30$ con $2$ d.o.f.
Nel goodness of fit, il confronto con il modello saturato d\`a $\chi^2_{\min}$ con $N-K$ d.o.f.

\textbf{(c)} Il frequentista non ha una PDF dei parametri, quindi non pu\`o integrare $\nu$. Usa il profile likelihood ratio
\[
\lambda(\mu)=\frac{L(\mu,\hat{\hat\nu}(\mu))}{L(\hat\mu,\hat\nu)},
\]
ovvero a ogni $\mu$ rimassimizza su $\nu$.
Il bayesiano integra, $p(\mu\mid x)=\int p(\mu,\nu\mid x)\,d\nu$, senza rapporti n\'e Wilks.
La marginalizzazione dipende dal prior su $\nu$ e tiene conto del volume dello spazio dei parametri; il profiling no.
Con un MCMC la marginale si ottiene ignorando le altre coordinate dei campioni.
}

\section{Limiti superiori, goodness of fit, combinazioni e intervalli credibili (Week 2--4)}

\Dom{(a)~Una singola misura $x_{\rm obs}$ di una Gaussiana con $\mu$ ignoto e $\sigma$ noto: come calcoli il limite superiore frequentista al livello $1-\alpha$? E quello bayesiano?
(b)~In un counting experiment con $n\sim\mathrm{Poisson}(s+b)$, come calcoli il $p$-value di scoperta e la significativit\`a? Come imposti un limite su $s$?
(c)~Come si fa un test di goodness of fit con il $\chi^2$? Cosa c'entrano il modello saturato e Wilks? Quando il $\chi^2$ non va bene?
(d)~Come combini pi\`u misure dello stesso parametro?
(e)~Intervalli credibili: equal-tailed contro HPD. Come si riassume un posterior in un punto?}

\Ris{
\textbf{(a)} Frequentista: $\mu_{\rm up}$ \`e il valore ipotetico di $\mu$ tale che, sotto quell'ipotesi, la probabilit\`a di osservare $x<x_{\rm obs}$ \`e $\alpha$; quindi $\mu_{\rm up}=x_{\rm obs}+\Phi^{-1}(1-\alpha)\,\sigma$ (per il $95\%$: $x_{\rm obs}+1.645\,\sigma$). \`E l'inversione di un test di ipotesi.
Bayesiano: con prior piatto $p(\mu\mid x)\propto L(\mu)$ e il limite al $95\%$ \`e il quantile $0.95$ del posterior; per una Gaussiana illimitata i numeri coincidono con quelli frequentisti, ma il significato \`e diverso (probabilit\`a su $\mu$ contro copertura su esperimenti ripetuti).
Con un vincolo fisico $\mu\ge0$ il bayesiano tronca il prior e il risultato dipende da esso; il frequentista pu\`o trovarsi con intervalli vuoti o poco informativi se $x_{\rm obs}$ \`e molto negativo. (Soluzioni come Feldman--Cousins e CLs sono in Cowan, non nelle slide.)

\textbf{(b)} Sotto $H_0$ ($s=0$) il numero atteso \`e $b$; il $p$-value di scoperta \`e
\[
p=P(n\ge n_{\rm obs}\mid b)=\sum_{n\ge n_{\rm obs}}\frac{b^ne^{-b}}{n!},
\]
e la significativit\`a $Z=\Phi^{-1}(1-p)$ (la scoperta richiede $Z>5$).
Limite superiore su $s$: il valore $s_{\rm up}$ tale che $P(n\le n_{\rm obs}\mid s_{\rm up}+b)=\alpha$ (stessa logica del punto (a), con una Poisson).
La significativit\`a dipende anche da quanto bene \`e noto $b$ (un parametro di disturbo da profilare o marginalizzare).

\textbf{(c)} Con $N$ misure $y_i$ e modello $f(x;\theta)$ a $K$ parametri, si confronta con il modello saturato ($f_{\rm sat}(x_i)=y_i$). Con errori gaussiani il rapporto di verosimiglianza diventa
\[
t=\sum_{i=1}^N\frac{\bigl(y_i-f(x_i;\hat\theta)\bigr)^2}{\sigma_i^2}\equiv\chi^2_{\min}.
\]
Per Wilks $t\sim\chi^2_n$ con $n=N-K$ (parametri del saturato meno quelli del modello). Si rifiuta il modello se il $p$-value $=1-F_{\chi^2_n}(t_{\rm obs})$ \`e piccolo (regione critica per $\chi^2$ grande); un $\chi^2/n\simeq1$ \`e atteso per un buon fit.
Il $\chi^2$ presume errori gaussiani; con pochi conteggi per bin vale meglio il rapporto di verosimiglianza di Poisson. Wilks vale nel limite di grandi campioni e con condizioni di regolarit\`a.

\textbf{(d)} La pratica migliore \`e combinare le likelihood, $\ln L(\theta)=\sum_i\ln L_i(\theta)$, non gli intervalli. Per questo conviene pubblicare la profile likelihood (per esempio su HEPData); se si ha solo l'intervallo, si pu\`o ricostruire la likelihood in modo approssimato. Gli esempi del corso sono le medie HFLAV ($R(D)$, $R(D^*)$) e i contorni al $68\%$ con la correlazione $\rho$ e il $p$-value del $\chi^2$ della combinazione.

\textbf{(e)} Intervallo credibile al livello CL: la regione che contiene il CL del posterior.
\emph{Equal-tailed}: lascia la stessa frazione di probabilit\`a nelle due code; \emph{HPD} (Highest Posterior Density): l'insieme dei $\theta$ con densit\`a maggiore di una soglia, \`e l'intervallo pi\`u corto, ma dipende dalla parametrizzazione.
In pi\`u parametri non c'\`e ordinamento, quindi non si definiscono code uguali: si usa l'HPD.
Punto: moda (MAP), mediana o media del posterior, scelta arbitraria che dipende dalla forma; con prior piatto il MAP coincide con l'MLE.
Significato: un intervallo credibile al $95\%$ dice ``dati questi dati e questo prior c'\`e il $95\%$ di probabilit\`a che il parametro sia qui''; un intervallo di confidenza dice ``il $95\%$ degli intervalli costruiti cos\`i in esperimenti ripetuti contiene il vero valore''.
}

\section{Bayes, MCMC e Bayes factor (Week 4)}

\Dom{(a)~Perch\'e $\pi(\theta)=\text{const}$ non rappresenta ignoranza (pensa a $\lambda=1/\theta$)? Quando MAP = MLE? Cosa \`e un prior improprio e quando va bene?
(b)~Descrivi Metropolis--Hastings. Cosa sono burn-in, autocorrelazione e convergenza? Perch\'e serve l'MCMC invece di una griglia?
(c)~Cos'\`e un Bayes factor e come si calcola l'evidenza? Perch\'e i prior impropri non vanno bene per il Bayes factor e dove si vede il rasoio di Ockham?}

\Ris{
\textbf{(a)} Con $\lambda=1/\theta$: $\pi(\lambda)=\pi(\theta)\,|d\theta/d\lambda|\propto1/\lambda^2$, non piatto. Il ``piatto'' dipende dalla parametrizzazione, e un'ignoranza completa non \`e ben definita.
Prior non informativi sensati riflettono le simmetrie: uniforme per un parametro di traslazione, $1/\theta$ per un parametro di scala (uniforme in $\ln\theta$); altre prescrizioni: Jeffreys, coniugati, massima entropia.
Con prior piatto $p(\theta\mid x)\propto L(\theta)$, quindi la moda del posterior (MAP) coincide con l'MLE.
Un prior improprio non \`e normalizzabile; va bene per stimare i parametri finch\'e il posterior \`e normalizzabile.
Va sempre controllata la sensibilit\`a del posterior al prior; se si pubblica la likelihood il risultato resta valido (``se questo prior, allora questo posterior'').

\textbf{(b)} Si parte da $\vec\theta_0$ e si propone $\vec\theta'$ da una distribuzione simmetrica $q$ (per esempio una Gaussiana centrata su $\vec\theta_0$). Se $p(\theta')>p(\theta_0)$ si accetta; altrimenti si accetta con probabilit\`a $p(\theta')/p(\theta_0)$ (si estrae $u\in[0,1]$ e si accetta se $u\le p(\theta')/p(\theta_0)$). Se si rifiuta, si ripete il vecchio punto.
Perch\'e serve: una griglia o l'accept--reject scalano esponenzialmente ($100$ punti $\times$ $5$ parametri $=10^{10}$), l'MCMC campiona soprattutto le regioni ad alta densit\`a e scala circa linearmente.
Burn-in: la catena impiega tempo per arrivare nella regione ad alta densit\`a, quindi si scartano i primi passi.
Autocorrelazione: punti consecutivi sono correlati, il numero effettivo di campioni \`e minore di $N$.
Convergenza: la catena pu\`o restare intrappolata in un massimo locale; si lanciano pi\`u catene da punti diversi.
Dal campione si ottengono le marginali dei POI ignorando le altre coordinate, e gli intervalli credibili (equal-tailed o HPD).

\textbf{(c)} Posterior odds $=$ Bayes factor $\times$ prior odds:
\[
\frac{P(H_1\mid x)}{P(H_0\mid x)}=\frac{P(x\mid H_1)}{P(x\mid H_0)}\cdot\frac{\pi(H_1)}{\pi(H_0)} ,\qquad K=\frac{P(x\mid H_1)}{P(x\mid H_0)}
\]
(si legge con la scala di Jeffreys). Ogni $P(x\mid H_i)$ \`e l'evidenza (marginal likelihood):
\[
P(x\mid H_i)=\int P(x\mid\theta_i,H_i)\,\pi_i(\theta_i\mid H_i)\,d\theta_i .
\]
Con un prior improprio l'integrale non \`e definito (la costante di normalizzazione \`e arbitraria, e con essa $K$).
Ockham: un modello con pi\`u parametri o pi\`u volume di prior diluisce la probabilit\`a e ha un'evidenza minore, quindi \`e penalizzato se i dati non ne hanno bisogno; la penalit\`a \`e incorporata nel metodo.
}
\Trap{Confondere il Bayes factor con un $p$-value: il $p$-value non confronta due modelli ed \`e frequentista.}

\section{Statistical learning, loss e regolarizzazione (L9, L10)}

\Dom{(a)~Scrivi la decomposizione dell'MSE di test e spiega cosa succede all'aumentare della flessibilit\`a. Cos'\`e l'overfitting?
(b)~Mostra che MSE corrisponde a una likelihood gaussiana e BCE a una di Bernoulli. Perch\'e la BCE \`e pi\`u aggressiva dell'MSE nella classificazione?
(c)~Differenza tra Ridge e LASSO; perch\'e LASSO azzera i coefficienti? Come si sceglie $\lambda$?}

\Ris{
\textbf{(a)}
\[
E\bigl[(Y-\hat f(x))^2\bigr]=\mathrm{Bias}[\hat f]^2+\mathrm{Var}[\hat f]+\sigma^2_\epsilon .
\]
Il bias \`e l'errore dovuto a un modello troppo semplice; la varianza dice quanto cambia $\hat f$ cambiando il training set; $\sigma^2_\epsilon$ \`e irriducibile (risoluzione del rivelatore, pile-up).
Con pi\`u flessibilit\`a il bias scende e la varianza sale: l'MSE di test \`e a U, quello di training scende sempre.
Overfitting: il modello segue il rumore del training (errore di training basso, di test alto, alta varianza); underfitting: alto bias.
Parametrici: veloci e interpretabili ma rigidi (rischio di alto bias); non parametrici: flessibili ma richiedono pi\`u dati (rischio di alta varianza). Inferenza $\Rightarrow$ interpretabilit\`a, predizione $\Rightarrow$ flessibilit\`a.

\textbf{(b)} Con $y=f(x)+\epsilon$ e $\epsilon$ gaussiano, $-\ln L=\text{const}+\sum_i (y_i-f(x_i))^2/(2\sigma^2)$: con $\sigma$ costante minimizzare l'NLL equivale a minimizzare l'MSE (il fit $\chi^2$ \`e lo stesso caso).
Con $y\in\{0,1\}$ Bernoulli e $p=f(x)$:
\[
-\ln L=-\sum_i\bigl[y_i\ln p_i+(1-y_i)\ln(1-p_i)\bigr]\quad(\text{BCE}).
\]
La BCE \`e pi\`u aggressiva perch\'e $-\ln p$ diverge quando si assegna probabilit\`a $\sim0$ alla classe vera, mentre l'MSE \`e limitato; punisce di pi\`u gli errori confidenti.
Con Ridge/LASSO in classificazione si usa la regressione logistica con penalit\`a (BCE con sigmoide), non la RSS.

\textbf{(c)} $\mathrm{Loss}_{\rm reg}=\mathrm{Loss}+\lambda R(\beta)$, con $R=\sum\beta_j^2$ (Ridge) o $\sum|\beta_j|$ (LASSO).
Ridge restringe i coefficienti ma non li azzera; LASSO fa feature selection. LASSO conviene con poche feature rilevanti, Ridge con molte di importanza simile.
Forma equivalente: si minimizza la RSS sotto un vincolo, $\sum\beta_j^2\le s$ (sfera, liscia) o $\sum|\beta_j|\le s$ (rombo con gli angoli sugli assi). L'ellisse della RSS tocca il vincolo L1 spesso in un angolo, dove qualche $\beta_j=0$.
Aumentando $\lambda$ cala la flessibilit\`a: meno varianza, pi\`u bias.
$\lambda$ si sceglie con la cross-validation ($k$-fold: si allena su $k-1$ fold e si valuta sul rimanente, si media, si prende il $\lambda$ con MSE di test minimo).
Prima serve lo standard scaler (media $0$, $\sigma\,1$), altrimenti la penalit\`a colpisce di pi\`u le feature con unit\`a piccole.
}

\section{PCA, clustering, alberi, bagging e boosting (L10, L11)}

\Dom{(a)~Come si definisce la prima componente principale? Relazione con SVD e autovalori; come scegli quante PC tenere? Differenza PCA/LASSO?
(b)~Descrivi l'algoritmo di Lloyd per il K-means e cosa pu\`o andare storto.
(c)~Come cresce un albero? Perch\'e Gini/entropia e non l'errore $0/1$? Cos'\`e il pruning? Bagging contro boosting: quale riduce varianza e quale bias?}

\Ris{
\textbf{(a)} La prima PC $\varphi_1$ (con $\|\varphi_1\|=1$) \`e la direzione lungo cui la varianza dei dati centrati \`e massima; equivale alla retta che minimizza la distanza quadratica dai punti. Le successive hanno la stessa definizione ma sono ortogonali alle precedenti.
Si ottengono dalla SVD $X=U\Sigma\Phi^T$: le colonne di $\Phi$ sono le PC, quelle di $U\Sigma$ gli score; equivalentemente sono gli autovettori di $X^TX$ (la covarianza), con varianza lungo la PC pari all'autovalore.
La varianza totale si conserva; $\mathrm{PVE}_m=\lambda_m/\sum\lambda$, e si tengono le PC fino a una frazione sufficiente della varianza (scree plot).
La PCA \`e non supervisionata e non usa il target: tiene le direzioni di massima varianza, non quelle utili per predire $Y$. LASSO \`e supervisionato e seleziona feature originali; la PCA combina tutte le feature in nuove direzioni. Va fatta dopo lo scaling.
Autoencoder: la PCA trova un iperpiano, l'AE una variet\`a non lineare.

\textbf{(b)} Si assegna a ogni punto un'etichetta a caso, poi si itera: (1)~calcola il centroide di ogni cluster, (2)~riassegna ogni punto al centroide pi\`u vicino, fino a stabilit\`a.
Minimizza la dissimilarit\`a intra-cluster $W(C_\ell)$. Termina sempre ma pu\`o fermarsi in un minimo locale (si lancia con pi\`u inizializzazioni).
$K$ va fissato a priori e dipende dall'applicazione; vanno controllate robustezza (sottocampioni, altri algoritmi) e scaling (senza scaling le variabili con varianza maggiore dominano la distanza euclidea). Ogni punto va in un solo cluster.

\textbf{(c)} Alberi: partizioni binarie ricorsive del feature space in rettangoli (tagli allineati agli assi), modelli non lineari e molto interpretabili.
Crescita greedy: a ogni nodo si scandiscono tutte le feature e tutti i valori di taglio, si sceglie quello che minimizza l'impurit\`a pesata dei figli. Non c'\`e lookahead, quindi si perde l'ottimo globale.
Impurit\`a: Gini $=\sum p(1-p)$ o entropia (information gain). L'errore $0/1$ \`e discontinuo e degenere (un taglio utile pu\`o non migliorarlo subito); Gini ed entropia sono continue e sensibili ai miglioramenti parziali.
Crescere fino a nodi puri \`e overfitting: si usano stopping (profondit\`a massima, dimensione minima, minima riduzione di impurit\`a) con CV, oppure il pruning a costo-complessit\`a, $\text{errore}+\alpha\cdot(\#\text{nodi})$, con $\alpha$ scelto in CV.
Regressione: loss MSE, predizione della foglia $=$ media dei $y$, si sceglie il taglio che minimizza la varianza.
\textbf{Bagging} (bootstrap $+$ media o voto): riduce la varianza di modelli a basso bias (alberi profondi).
\textbf{Boosting}: modelli molto regolarizzati addestrati in sequenza, con pi\`u peso agli esempi sbagliati; riduce il bias di modelli a bassa varianza (AdaBoost, gradient boosting, XGBoost; parametri: learning rate, numero di alberi, subsample).
ROC: TPR contro FPR al variare della soglia; l'AUC la riassume.
}

\section{Reti neurali, backpropagation e ottimizzazione (L12, L13)}

\Dom{(a)~Cosa calcola un neurone, perch\'e servono non linearit\`a? Cosa dice e cosa non dice l'approssimazione universale? Attivazioni per hidden e output.
(b)~Discesa del gradiente, SGD con mini-batch, backprop, vanishing gradient, batch normalization.
(c)~Momentum, RMSProp, Adam; dropout, early stopping, data augmentation; dataset sbilanciato.}

\Ris{
\textbf{(a)} Un neurone calcola $a=\sigma(w\cdot x+b)$. Senza non linearit\`a una composizione di strati lineari resta lineare; con la non linearit\`a si imparano relazioni tra le feature (esempio XOR).
Approssimazione universale: una rete con uno strato nascosto, finito e con attivazione non lineare, approssima qualunque funzione continua su un compatto. Non dice quanti neuroni servono n\'e come trovare i pesi: tutta la pratica del deep learning \`e questo.
Perch\'e la profondit\`a: servono meno parametri, le sottofeature si riutilizzano, la sovraparametrizzazione sembra rendere la loss pi\`u liscia.
Hidden: ReLU (non satura, nessun vanishing gradient per $x>0$); la sigmoide \`e poco usata (gradiente piccolo, output non centrato).
Output: lineare per la regressione, sigmoide per la binaria, softmax per la multiclasse mutuamente esclusiva (somma $1$). Loss: MSE, BCE, cross-entropy categorica.

\textbf{(b)} $w\leftarrow w-\eta\nabla_wL(w)$.
GD usa tutto il dataset a ogni passo (troppo costoso); SGD un esempio, mini-batch un piccolo gruppo ($8$--$32$): il gradiente \`e stimatore non distorto di quello vero, con rumore utile, e sfrutta GPU e parallelismo.
Backprop: la regola della catena $\partial L/\partial w=\partial L/\partial\hat y\cdot\partial\hat y/\partial a\cdot\partial a/\partial w$ propagata all'indietro strato per strato; dice ``quanta colpa'' ha ogni peso nell'errore.
Vanishing gradient: il prodotto di molte derivate piccole (sigmoide: massimo $0.25$) fa svanire il gradiente, i primi strati non imparano. Rimedi: ReLU, inizializzazione, normalizzazione, skip connections.
Batch normalization: normalizza le attivazioni per mini-batch e stabilizza l'addestramento (motivazione originale: internal covariate shift).
La loss \`e non convessa: dipende dall'inizializzazione, ha minimi locali e selle (frequenti in alta dimensione); i minimi piatti generalizzano meglio di quelli stretti. Il learning rate annealing aiuta.

\textbf{(c)} \textbf{Momentum}: accumula una velocit\`a, smorza le oscillazioni nelle direzioni ad alta curvatura e accelera dove il gradiente \`e costante; aiuta con le selle.
\textbf{AdaGrad/RMSProp}: learning rate per parametro, scalato con la somma (AdaGrad) o la media mobile (RMSProp) dei gradienti al quadrato.
\textbf{Adam}: momentum (primo momento $m_t$) $+$ scala adattiva (secondo momento $v_t$), con correzione del bias iniziale; $\beta_1=0.9$, $\beta_2=0.999$, $\epsilon=10^{-8}$.
\textbf{Dropout}: spegne a caso una frazione dei neuroni in training, impedisce di affidarsi a poche feature.
\textbf{Early stopping}: ci si ferma quando la loss di validazione smette di scendere. \textbf{Noise injection}: rumore su input o etichette. \textbf{Data augmentation}: trasformazioni che preservano la fisica.
Dataset sbilanciato: oversampling della minoranza (rischio overfitting); undersampling della maggioranza (si perde informazione); SMOTE (punti sintetici per interpolazione tra vicini della minoranza: meno overfitting, possibile bias); class weights (downsampling $+$ upweighting, cos\`i il modello impara sia le feature sia la vera distribuzione delle classi).
Preprocessing: standard scaler, e min-max spesso per le reti.
}

\section{Immagini, filtri e preparazione dei dati (L14, L10, L13)}

\Dom{(a)~Cos'\`e un'immagine digitale e cosa mostra un istogramma? Che differenza c'\`e tra operazioni puntuali e filtri? Che cosa fanno thresholding, power law ed equalizzazione?
(b)~Cos'\`e una convoluzione, perch\'e \`e importante e come si gestiscono i bordi? Cosa fanno i filtri di smoothing e quelli derivativi?
(c)~``Garbage in, garbage out'': quali problemi hanno i dati reali e come si prepara un dataset? Standard scaler contro min-max; learning curve e iperparametri.}

\Ris{
\textbf{(a)} Un'immagine \`e una matrice 2D di intensit\`a (livelli di grigio $0$--$255$, oppure $3$ canali RGB, quindi larghezza$\times$altezza$\times3$). La digitalizzazione \`e un'approssimazione della scena reale (risoluzione spaziale e livelli di intensit\`a).
L'istogramma conta quante volte compare ogni intensit\`a: mostra luminosit\`a, contrasto, esposizione e dynamic range; la cumulata serve per l'equalizzazione.
\textbf{Operazioni puntuali}: la nuova intensit\`a dipende solo da quella vecchia e dalla funzione di mappatura, non dalla posizione n\'e dai vicini. Esempi: \emph{thresholding} (da grigio a binario), \emph{power law} (comprime o espande intervalli scuri o chiari, mettendo in evidenza dettagli diversi), \emph{equalizzazione} (rende l'istogramma il pi\`u uniforme possibile, usando tutto il range dinamico).
\textbf{Filtri}: combinano il pixel con i suoi vicini (dimensione, forma e pesi della finestra; lineari o non lineari).

\textbf{(b)} Un filtro lineare e invariante per traslazione \`e una convoluzione (teorema): ogni pixel diventa la somma pesata del suo intorno con il kernel.
Bordi: crop, pad, extend, wrap.
Smoothing: la media su una finestra (blurring) o il filtro gaussiano, con pesi maggiori al centro e normalizzazione preservata; $\sigma$ \`e la larghezza.
Edge detection: i bordi sono salti di luminosit\`a, individuabili dalla derivata prima; in discreto si usa un kernel centrato e simmetrico come $[-0.5,0,0.5]$.
Una CNN usa gli stessi filtri ma li impara dai dati, invece di fissarli a mano.

\textbf{(c)} La qualit\`a dell'output dipende da quella dell'input. Tipici problemi: dati inaccurati, mancanti, distorti (non rappresentativi), rumorosi, outlier. Rimedi: data cleaning, feature engineering, validazione; \`e spesso la parte che richiede pi\`u tempo.
\textbf{Standard scaler}: media $0$ e $\sigma\,1$ per ogni feature; evita che una feature domini solo per le unit\`a e aiuta la convergenza. \textbf{Min-max}: comprime in un intervallo, di solito $[0,1]$, spesso preferito per le reti.
Lo scaler va adattato solo sul training set e poi applicato a validazione e test, per non far entrare informazione dal test nel training.
Learning curve: l'errore di training e di validazione in funzione delle epoche; il plateau dipende da dimensione del training set, architettura e inizializzazione; se l'errore di validazione risale c'\`e overfitting.
Gli iperparametri non si apprendono con la discesa del gradiente ma con euristiche (prove manuali, grid search, random search, AutoML), valutate su un set di validazione; il test set si usa una sola volta alla fine.
Il lavoro standard: baseline semplice (modello lineare) per riferimento, poi BDT, poi rete, confrontando con la stessa metrica.
}

\section{CNN, RNN/LSTM e GNN (L14--L16)}

\Dom{(a)~Perch\'e una FFNN su immagini non scala? Le due assunzioni di una CNN? Dimensione dell'output di una convoluzione. Cosa fa il pooling?
(b)~Perch\'e le RNN hanno vanishing/exploding gradient e perch\'e la tanh? Gate di un LSTM e perch\'e la cell state aiuta. Windowing e forecasting ricorsivo contro diretto.
(c)~Rappresentazione di un grafo, message passing, invarianza/equivarianza, over-smoothing.}

\Ris{
\textbf{(a)} Una FFNN appiattisce l'immagine e perde le relazioni spaziali: con $64\times64$ e $500$ nodi nascosti servono $\sim2$ milioni di pesi (il triplo con il colore).
Assunzioni della CNN: \textbf{localit\`a} (pattern vicini contano di pi\`u) e \textbf{invarianza per traslazione} (stessi pesi in ogni posizione). I filtri sono appresi, i parametri calano molto, le operazioni si parallelizzano sulla GPU.
Output: $o=\left\lfloor\dfrac{n+2p-m}{s}\right\rfloor+1$ ($n$ input, $m$ kernel, $p$ padding, $s$ stride); il padding ``same'' mantiene la dimensione.
Con pi\`u canali (RGB) ogni canale ha il suo filtro e i risultati si sommano; si aumenta il numero di filtri negli strati successivi.
Pooling (max o average): riduce la dimensione, introduce invarianza locale, riduce l'overfitting.
Per immagini sparse (jet) si usano filtri pi\`u grandi (es.\ $11\times11$). Dropout e data augmentation valgono come per le DNN.

\textbf{(b)} Una RNN legge la sequenza e aggiorna uno stato $h_t=\tanh(Wh_{t-1}+Ux_t+b)$ con pesi condivisi nel tempo; accetta sequenze di lunghezza variabile.
La backpropagation through time moltiplica molte derivate: piccole $\Rightarrow$ il gradiente svanisce; pesi grandi $\Rightarrow$ esplode.
La tanh \`e preferita: output in $[-1,1]$ (pu\`o sommare e sottrarre informazione), centrata in zero, derivata massima $1$ (contro $0.25$ della sigmoide); ReLU farebbe esplodere i gradienti.
Per dipendenze lunghe una RNN semplice sovrascrive la memoria e il gradiente va a zero.
\textbf{LSTM}: cell state $C_t$; forget gate ($0$ dimentica, $1$ mantiene); input gate (cosa aggiornare e candidati nuovi); aggiornamento $C_t=f\cdot C_{t-1}+i\cdot\tilde C_t$; output gate (cosa mandare fuori). Grazie alla somma il gradiente scorre all'indietro senza svanire facilmente.
Windowing: la finestra di input ha una larghezza, un offset rispetto alla label e un insieme di feature. \code{return\_sequences=True} d\`a un output per passo (per impilare LSTM), \code{False} solo l'ultimo.
Servono stazionariet\`a (debole: media, varianza e autocorrelazione costanti), feature engineering (angoli del vento, orario) e una baseline semplice (ripetere l'ultimo valore).
Ricorsivo: la predizione diventa input (rischio di feedback e bias); diretto: si predice direttamente l'orizzonte voluto.

\textbf{(c)} Un grafo \`e $G=(V,E)$ con matrice di adiacenza $A$ ($|V|\times|V|$, simmetrica se non orientato), matrice $X$ delle feature dei nodi ($|V|\times D$) e, opzionalmente, delle feature degli archi. $A^m$ conta i cammini di lunghezza $m$.
Message passing: \textbf{aggregazione} (somma degli embedding dei vicini) e \textbf{aggiornamento} (trasformazione lineare del nodo e del messaggio, poi attivazione non lineare). Dopo $K$ strati ogni nodo contiene informazione sul suo intorno.
Compiti: classificazione di grafo (global pooling, ad esempio media, poi MLP), di nodo, di arco (link prediction).
Invarianza: l'ordine dei nodi \`e arbitrario, quindi $f(PV)=f(V)$ per l'output di grafo; l'output di nodo \`e equivariante: $f(PV,PAP^T)=Pf(V,A)$. Cos\`i la rete impara relazioni, non posizioni nella lista.
Bias induttivo relazionale: favorisce le soluzioni coerenti con la struttura (se sbagliato pu\`o essere troppo restrittivo).
Batching: grafi di dimensioni diverse come componenti disgiunte di un unico grafo grande, con pooling per grafo.
\textbf{Over-smoothing}: dopo troppe iterazioni gli embedding convergono allo stesso valore. Si scelgono gli strati vicino al diametro del grafo e si usano skip connection o MLP profondi per il passaggio dei messaggi.
Tipici HEP: jet tagging, tracking (classificazione degli archi); vantaggio: dati sparsi e irregolari, senza griglia.
}

\section{Autoencoder, VAE e incertezza in ML (L17, L18)}

\Dom{(a)~Come funziona un autoencoder e in cosa differisce dalla PCA? Anomaly detection e denoising; perch\'e serve il collo di bottiglia?
(b)~Prior, likelihood e posterior variazionale in un VAE; ELBO e i suoi termini; reparametrization trick; cosa succede con KL $=0$?
(c)~Incertezza aleatorica ed epistemica; bagging, deep ensemble, MC dropout, Bayes by Backprop; perch\'e una rete $(\mu,\sigma)$ cattura solo l'aleatorica?}

\Ris{
\textbf{(a)} Una FFNN con input e output della stessa forma: encoder $\to$ collo di bottiglia $\to$ decoder, che minimizza l'errore di ricostruzione (MSE). \`E addestramento self-supervised: nessuna etichetta.
Rispetto alla PCA impara una variet\`a non lineare a bassa dimensione, la PCA un iperpiano; in entrambi i casi si riduce la dimensione.
Il collo di bottiglia (undercomplete) obbliga a comprimere; senza, la rete impara l'identit\`a e non estrae nulla.
Anomaly detection: si addestra solo sul fondo, che ricostruisce bene (MSE basso); un evento anomalo, come un possibile segnale, ricostruisce male (MSE alto). \`E una classificazione a una classe.
Denoising: input corrotto (rumore o mascheramento), target pulito; l'AE impara le feature rilevanti e scarta il rumore.

\textbf{(b)} L'encoder produce una distribuzione $q_\phi(z\mid x)$ ($\mu$ e $\sigma$); il decoder \`e $p_\theta(x\mid z)$ (likelihood); il prior $p(z)$ \`e tipicamente $\mathcal N(0,I)$.
Il posterior vero $p(z\mid x)$ \`e intrattabile (richiede l'integrale su $z$) e si approssima con $q_\phi$ minimizzando la KL.
\[
\mathrm{ELBO}=E_{q}\bigl[\ln p_\theta(x\mid z)\bigr]-\mathrm{KL}\bigl(q_\phi(z\mid x)\,\|\,p(z)\bigr);\qquad\text{loss}=-\mathrm{ELBO}.
\]
Il primo termine \`e la ricostruzione (accuratezza fisica); il secondo regolarizza lo spazio latente verso il prior (continuo e regolare).
Massimizzare l'ELBO equivale a massimizzare un limite inferiore dell'evidenza $\ln p(x)$: si trasforma un integrale in un problema di ottimizzazione.
Reparametrization trick: $z=\mu+\sigma\epsilon$ con $\epsilon\sim\mathcal N(0,1)$; il campionamento diventa deterministico rispetto ai parametri e il gradiente passa (non si pu\`o derivare un nodo casuale).
Con KL $=0$ (o $\epsilon=0$) si torna a un AE deterministico, senza vincolo sul posterior.
Dopo l'addestramento il decoder \`e un generatore: si campiona $z\sim\mathcal N(0,I)$ e si ottengono eventi nuovi. Lo spazio latente regolare \`e interpretabile; la distanza dal centro del prior serve per le anomalie.
Banda di incertezza: si campiona da $\mathcal N(\mu,\sigma)$ per ogni input e si passano i campioni nel decoder.

\textbf{(c)} \textbf{Aleatorica}: rumore intrinseco nelle osservazioni (omoschedastica se costante, eteroschedastica se dipende dall'input); non si riduce con pi\`u dati.
\textbf{Epistemica}: ignoranza sui parametri del modello; si riduce con pi\`u dati ed \`e cruciale fuori distribuzione (open set recognition: una softmax assegna comunque un'alta probabilit\`a).
\textbf{Bagging}: $M$ reti su bootstrap; la media \`e la predizione, la varianza tra le reti \`e l'incertezza epistemica.
\textbf{Deep ensemble}: stessa architettura, inizializzazioni diverse, stesso dataset; semplice e parallelizzabile, ma costoso.
\textbf{Rete $(\mu,\sigma^2)$} con NLL gaussiana: cattura il rumore dei dati (aleatorica); \`e deterministica e non segnala di essere fuori distribuzione, quindi manca l'epistemica. Versione pi\`u flessibile: quantile regression.
\textbf{MC dropout}: dropout attivo anche in test e $T$ passaggi; equivale a un'inferenza variazionale approssimata (Gal), la varianza tra i passaggi misura l'ignoranza.
\textbf{Bayes by Backprop}: ogni peso \`e una distribuzione $\mathcal N(\mu_i,\sigma_i)$ (parametri $\mu_i,\rho_i$: $2\times$ i parametri), con reparametrization trick; la loss \`e l'ELBO (likelihood $+$ KL con il prior, costo di complessit\`a). Con pi\`u dati la likelihood domina sul prior. Pi\`u sensibile agli iperparametri.
VAE contro BNN: stesso motore (ELBO e reparametrization); il VAE \`e stocastico solo nel collo di bottiglia (incertezza nella rappresentazione, generazione), la BNN in tutti i pesi (incertezza del modello).
}

\section{Unfolding e OmniFold (L19, L20)}

\Dom{(a)~Equazione di folding discreta; cos'\`e $R$ e cosa dice la somma per colonna? Perch\'e invertire $R$ \`e mal posto? Cosa dice la covarianza tra bin vicine?
(b)~Confronta inversione di matrice, bin-by-bin, Tikhonov e IBU di D'Agostini. Cos'\`e il forward folding?
(c)~Come funziona OmniFold? Trucco del classificatore, le due fasi di reweighting, perch\'e \`e unbinned e multidimensionale; perch\'e la curse of dimensionality uccide l'IBU?}

\Ris{
\textbf{(a)} Il numero atteso di eventi misurati nella bin $i$ \`e $\mu_i=\sum_j R_{ij}\nu_j+\beta_i$, con $\nu$ i veri conteggi per bin e $\beta$ il fondo.
$R_{ij}=P(\text{misurato in }i\mid\text{generato in }j)$ si costruisce con l'MC (risoluzione, accettanza, efficienza); la somma su $i$ per una colonna $j$ d\`a l'efficienza $\varepsilon_j$.
L'inversione $\hat\nu=R^{-1}(n-\beta)$ \`e l'MLE: azzera la distanza tra il modello ripiegato e i dati, quindi \`e un overfitting analitico ($N$ punti, $N$ parametri) con varianza enorme a bin piccole.
\`E un problema mal posto: $R$ smussa i dettagli fini e $R^{-1}$ li ricrea dalle fluttuazioni statistiche, trattandole come segnale; piccole variazioni dei dati danno grandi variazioni del risultato.
La covarianza $V=R^{-1}V_n(R^{-1})^T$ mostra anticorrelazione tra bin adiacenti (una fluttuazione in una va sottratta ai vicini per conservare il totale).
Bin larghe rispetto alla risoluzione: stabile ma con bias; bin piccole: poco bias ma instabile. \`E il bias--varianza di nuovo.
$R$ dipende dalla distribuzione vera usata nell'MC: il bias \`e nullo se la distribuzione \`e piatta dentro la bin.

\textbf{(b)} \textbf{Bin-by-bin}: fattore moltiplicativo per bin $C_i=\nu_i^{\rm MC}/\mu_i^{\rm MC}$; ipotizza migrazioni trascurabili ($R$ diagonale, purity $>70$--$80\%$) e forma dell'MC corretta; va bene con bin larghe e per un primo sguardo.
\textbf{Tikhonov}: massimizza $\ln L(\nu)+\tau S(\nu)$, con $S$ che penalizza la curvatura ($k=2$) o la pendenza ($k=1$); $\tau$ controlla il compromesso bias--varianza. \`E un prior sulla regolarit\`a: la smoothness deve essere informazione che hai davvero.
\textbf{IBU}: Bayes sulle cause $C_j$ e gli effetti $E_i$,
\[
P(C_j\mid E_i)=\frac{P(E_i\mid C_j)P(C_j)}{\sum_k P(E_i\mid C_k)P(C_k)},
\]
ogni evento misurato in $i$ \`e riassegnato alla bin vera $j$ con quel peso, dividendo per l'efficienza.
Il prior \`e la distribuzione vera che si cerca (argomento circolare), risolto iterando: si parte da $P_0$ (MC o piatto), si normalizza la stima in nuovo prior e si ripete.
Poche iterazioni: domina il prior (bias); troppe: si dimentica il prior e crescono le fluttuazioni (varianza); lo stop anticipato \`e la regolarizzazione (tipicamente $3$--$6$ iterazioni).
\textbf{Forward folding}: si confronta la teoria ripiegata con i dati (fit $\chi^2$ o likelihood), senza invertire nulla; spesso \`e la scelta migliore (``Don't unfold'').

\textbf{(c)} Con molte variabili la matrice di risposta ha dimensioni enormi ($5{+}5$ variabili con $20$ bin ciascuna: $\sim10^{13}$ bin), quindi l'IBU non scala.
\textbf{Trucco}: un classificatore binario addestrato con BCE per distinguere campioni da $p(x)$ e $q(x)$ d\`a $f/(1-f)=p/q$, il rapporto di verosimiglianza: da stima di densit\`a (difficile) a classificazione (facile). I pesi sono $w(x)=p(x)/q(x)$.
\textbf{Iterazione}:
\begin{enumerate}
\item reweighting della simulazione a livello reco verso i dati reali (primo classificatore);
\item i pesi si riportano sugli eventi a livello pre-detector (pull-back);
\item un secondo classificatore, tra eventi pre-detector e pre-detector ripesati, d\`a nuovi pesi (smoothing);
\item si ripete; alla fine gli eventi del generatore ripesati sono la stima.
\end{enumerate}
Il risultato \`e unbinned, event-based, likelihood-free e multidimensionale: sfrutta tutte le correlazioni e si possono unfoldare molte osservabili insieme.
Assunzione: la simulazione deve modellare bene la risposta del detector.
In sostanza OmniFold \`e l'IBU unbinned, multidimensionale e basato su reti neurali.

\textbf{(d) Progetto} \emph{(da completare)}: dove hai usato idee del corso nello studio di ACP su LHCb open data (likelihood fit, profile likelihood, classificatori, incertezze), e quali assunzioni e punti deboli ti aspetti che ti contestino. Questa risposta va scritta dopo aver riletto il report.
}

\section{Confronti e domande trabocchetto}

Domande brevi, in cui si cade spesso. Rispondi prima di leggere la risposta.

\begin{description}[style=nextline,leftmargin=0.5cm]
\item[``Il $p$-value \`e la probabilit\`a che $H_0$ sia vera?''] No. \`E $P(\text{dati almeno cos\`i estremi}\mid H_0)$. $P(H_0\mid\text{dati})$ richiede un prior (Bayes).
\item[``Un intervallo di confidenza al $68\%$ contiene $\mu$ con il $68\%$ di probabilit\`a?''] Per un frequentista no: $\mu$ \`e fisso, l'intervallo fluttua; il $68\%$ \`e la coverage su esperimenti ripetuti. Quella frase vale per l'intervallo credibile.
\item[``Prior piatto = ignoranza?''] No, dipende dalla parametrizzazione ($\lambda=1/\theta$ d\`a $1/\lambda^2$); servono prior che riflettano le simmetrie.
\item[``L'MLE \`e non distorto?''] Non in generale ($\hat\sigma^2$ con $1/n$); lo \`e solo asintoticamente. L'MLE trasformato (invarianza funzionale) \`e spesso distorto.
\item[Perch\'e Wilks pu\`o non valere?] Pochi dati, parametri al bordo dello spazio fisico (esempio $\mu\ge0$), likelihood non regolare: $2\Delta\mathrm{NLL}$ non segue un $\chi^2$ e va usato un MC (pseudo-esperimenti).
\item[Profiling contro marginalizzazione?] Profiling: massimizza su $\nu$ per ogni $\mu$ (frequentista, non dipende dal prior). Marginalizzazione: integra su $\nu$ con il suo prior (bayesiano, tiene conto del volume).
\item[Se NP \`e ottimo, perch\'e il ML?] Perch\'e le PDF non si conoscono in alta dimensione; il classificatore ottimo approssima il rapporto di verosimiglianza. \`E anche la base di OmniFold.
\item[Perch\'e MSE e BCE?] Sono NLL: MSE da rumore gaussiano, BCE da una Bernoulli. La loss \`e la likelihood cambiata di segno.
\item[Ridge contro LASSO?] L2 riduce, L1 azzera (angoli del vincolo). Ridge se molte feature contano, LASSO se ne bastano poche.
\item[Bagging contro boosting?] Bagging: riduce la varianza di modelli profondi (basso bias). Boosting: riduce il bias di modelli deboli e regolarizzati.
\item[Alto AUC e buon classificatore?] L'AUC non dice nulla sulla calibrazione n\'e sul punto di lavoro; con classi sbilanciate l'accuratezza inganna. Contano soglia, TPR/FPR, significance e il confronto con una baseline.
\item[Validazione o test?] La validazione serve a scegliere iperparametri, il test a stimare l'errore finale e si usa una volta.
\item[PCA contro AE contro VAE?] PCA: lineare, analitica (SVD), massima varianza. AE: non lineare, minimizza l'errore di ricostruzione. VAE: AE probabilistico con prior sul latente (KL) e ELBO, quindi generativo e con spazio regolare.
\item[Ensemble contro MC dropout contro Bayes by Backprop?] Ensemble: $M$ addestramenti, frequentista, costoso. MC dropout: un solo addestramento, equivalente a un'inferenza variazionale approssimata. BBB: pesi come distribuzioni con ELBO, $2\times$ parametri e pi\`u sensibile agli iperparametri.
\item[CNN contro GNN?] La CNN opera su una griglia regolare (invarianza per traslazione); la GNN su grafi irregolari e sparsi (invarianza per permutazioni dei nodi). Una convoluzione su immagine \`e un caso particolare di message passing.
\item[Perch\'e l'IBU non scala e OmniFold s\`i?] L'IBU bina in $d$ dimensioni (la matrice cresce come $(\text{bin})^{2d}$); OmniFold usa classificatori che lavorano direttamente sui singoli eventi (unbinned).
\item[Unfolding o forward folding?] Il forward folding (ripiegare la teoria e confrontare con i dati) \`e pi\`u semplice e senza bias da regolarizzazione; l'unfolding serve per confrontare esperimenti diversi e con teorie future.
\item[Perch\'e la regolarizzazione nell'unfolding introduce bias?] Perch\'e aggiunge un'informazione (smoothness) che stai assumendo; deve essere informazione vera.
\item[Frequentista contro bayesiano in una frase?] Frequentista: i dati sono casuali, il parametro \`e fisso (probabilit\`a = frequenza). Bayesiano: il dato \`e fisso, il parametro \`e una variabile casuale (probabilit\`a = credenza), con prior. Coincidono con molti dati e prior deboli; la likelihood \`e comune.
\end{description}

\end{document}
